\documentclass{article}
\usepackage{graphicx} % Required for inserting images
\usepackage{amssymb} % Required for math set font
\usepackage{float} % Required for inline figure placement
\usepackage[numbers]{natbib} % Required for numerical citations
\usepackage{mathtools} % Required for ceiling funciton notation
\DeclarePairedDelimiter{\ceil}{\lceil}{\rceil}

\title{Implementing Discrete Logarithm Algorithms}
\author{RORY ALEXANDER POULTER}
\date{October 2026}

\begin{document}

\maketitle

\section{Project Plan}

\begin{itemize}
    \item Project Supervisor: Peter Davies-Peck
    \item Project title: How group order smoothness and parallelism affect the optimal strategy for computing discrete logarithms.
    \item {Deliverables:
        \begin{itemize}
            \item Pollard's rho implementation (Basic)
            \item Pollard's kangaroo implementation (Basic)
            \item Test suite for algorithms with time and memory metrics (Basic)
            \item Test basic algorithms for use with Pohlig-Hellman subproblems (Basic)
            \item Generate scheme for choosing primes based on smoothness (Basic)
            \item Index calculus implementation (Intermediate)
            \item Parallel Pohlig-Hellman implementation (Basic)
            \item Parallel Index Calculus (Intermediate)
            \item Test smaller primes on local machine (Basic)
            \item Test large primes on NCC (Advanced)
        \end{itemize}}
\end{itemize}



\newpage
\section{Literature Survey}
\subsection{The Discrete Logarithm Problem}
\subsubsection{Notation}

Let $p$ be a prime and let $G = (\mathbb{Z}/p\mathbb{Z})^\times$ denote the multiplicative group of integers modulo $p$. The order $N$ of the group $G$ is the number of elements in $G$. The group contains some integer $g$ called a generator such that for all $h \in G$, $h$ can be obtained by the repeated multiplication of $g \mod p$. This can written as $h \equiv g^x \ (\mod  p)$, where $x \in \mathbb{Z}$ is referred to as the discrete logarithm. The subgroup $\langle g \rangle \subseteq G$ contains the set of numbers generated by $g$; this subgroup is finite (i.e. $| \langle g \rangle |$ is finite) and cyclic (i.e. $g^x \ (\mod p) = g^{x + iN} \ (\mod p) \ \forall i \in \mathbb{Z}$).

The process of finding this value $x$ where the values of $g, h, p$ are known is called the Discrete Logarithm Problem (DLP).

\subsubsection{Example Group}
Take $p = 11$ and $g = 2$. First check that $g$ generates the group by computing its powers mod $p$.

\begin{table}[H]
    \centering
    \begin{tabular}{c|c|c|c|c|c|c|c|c|c|c}
        $x$ & 0 & 1 & 2 & 3 & 4 & 5 & 6 & 7 & 8 & 9 \\
        $g^x \mod p$ & 1 & 2 & 4 & 8 & 5 & 10 & 9 & 7 & 3 & 6
    \end{tabular}
    \caption{The values generated by $g$ to the powers of 0 to 9}
    \label{tab:group-example}
\end{table}

Every non-zero residue 1 to 10 appears exactly once, so 2 is a generator and the group has order $N = p - 1 = 10$.

\subsection{Discrete Logarithm Algorithms} \label{sec:discrete-log-algorithms}

There are currently no known classical, efficient algorithms for solving discrete logarithms in all cases. The time complexity of different algorithms is measured against the bit length of the group order $b = log_2(N)$.

\subsubsection{Trial Multiplication (Brute Force)}

The most basic approach for solving the DLP is a brute force approach: test values $g^1, g^2, g^3, ...$ in order until the value is found. This has a time complexity of $O(2^b)$. Due to its inefficiency, it does not have practical applications.

\subsubsection{Shank's Algorithm}

Shank's Algorithm \cite{Shanks1971Class}, also known as Baby-step Giant-step, was developed in 1971 and is an optimisation of trial multiplication.

Let $x = im + j$ such that $m = \ceil {\sqrt{N}}, \ 0 \leq i,j < m$. Therefore $g^{im+j} = h \implies g^j = h(g^{-m})^i$. Pairs of values $(j, g^j)$ are calculated and incremented values of $i$ are tested until a match with a single value of $j$ is found.

It reduces the time complexity from $O(2^b)$ achieved with brute force to $O(2^{b/2})$ \cite{Menezes2001Handbook}.

\subsubsection{Pollard's Rho}

Pollard's Rho \cite{Pollard1978Monte} algorithm was developed in 1978 and is based on his algorithm for integer factorisation \cite{Pollard1975A}. The algorithm computes 4 integers $a, A, b, B$ such that $g^a h^b=g^Ah^B$. By replacing $h$ with $g^x$, you get the equation $(B-b)x = (a - A) \ (\mod N)$; solutions to this can be obtained using the extended Euclidean algorithm.

Similarly to Shank's algorithm, it has a time complexity of $O(2^{b/2})$ \cite{Menezes2001Handbook}.

\subsubsection{Pollard's Kangaroo}

Pollard's Kangaroo algorithm \cite{Pollard1978Monte}, also known as Pollard's Lambda algorithm, was also developed by John Pollard in the same 1978 paper as his Rho algorithm. It is probabilistic, and thus is not guaranteed to find the target. This is because it uses a pseudorandom map to generate steps. It searches for the value of $x$ in a range $[a, b] \subset \mathbb{Z}/p\mathbb{Z}$. The time complexity is $O(\sqrt{b - a})$ where $b$ in this case is the top of the range, and not the number of bits required to represent $N$.

\subsubsection{Pohlig-Hellman}

The Pohlig-Hellman algorithm \cite{Pohlig1978An} works by finding subgroups of $\langle g \rangle$. It does not solve for discrete logarithms directly, but rather decomposes the problem into simpler subproblems, which can then be solved with another algorithm, such as Shank's.

First, find the prime factorisation of $N = p_1^{n_1} \times p_2^{n_2} \times ... \times p_r^{n_r}$. Then let $r_i = N / p_i^{n_i}, \ g_i = g^{r_i}, \ h_i = h^{r_i} \ (\mod p), \ x_i = x \mod p_i^{n_i} = log_{g_i}h_i$. Once each $x_i$ is found, use Gauss' algorithm \cite{Gauss1986Disquisitiones} to find $x$. Because it utilises the prime factorisation of $N$, it works best when the group's order is a smooth integer.

If $N$ is smooth, the time complexity for the Pohlig-Hellman algorithm is $O(\sum_{i=1}^r n_i(log(n) + \sqrt{p_i}))$ \cite{Menezes2001Handbook}. 

The nature of the algorithm allows for easy parallelisation, which improves execution times \cite{Sakre2017Parallel}. A parallel implementation would use the following structure: the prime factorisation is processed on one core; each subgroup is given a single thread to calculate $x_i$; and finally the values of $x_i \forall i$ are used to calculate $x$.

\subsubsection{Index Calculus}

The index calculus algorithm \cite{Adleman1979A} is another discrete algorithm, but is different from the above algorithms. Where they only needed $(g, p, h)$ as inputs, the index calculus algorithm also requires a factor base as an input. This is a set of $r+1$ numbers $F_r = \{-1, 2, 3, ..., p_{r}\}$ containing -1 and the first $r$ primes. The algorithm is performed in three stages.

The first stage involves finding a set of $r$ linearly independent relations between the $F_r$ and the power of $g$. Each of these relation contributes an equation to a system of linear equations in $r$ unknowns: the discrete logarithms of the $r$ primes in $F_r$. The second stage solves these equations. This is often performed on a supercomputer as it demands large amounts of memory. The third stage searches for an integer $s$ such that $g^sh = (-1)^{f_0}2^{f_1}3^{f_2}...p_r^{f_r}$. Finally, rearrange it to solve $x=f_olog_g(-1) + f_1log_g2 + f_2log_g3 + ... + f_rlog_gp_r - s$.

As Adleman showed, the index calculus algorithm is subexponential, putting it far ahead of the algorithms shown above.

\subsubsection{General Number Field Sieve}

The General Number Field Sieve (GNFS) \cite{Lenstra1993The} is a variation of index calculus, differing in how it obtains smooth objects. Where index calculus finds integers $g^k \mod p$ that are smooth over small primes, GNFS finds pairs of values - an integer and an algebraic norm - both smooth.

A polynomial $f(x) = c_0+c_1x+c_2x^2+...+c_dx^d$ of degree $d$ is chosen, and an integer $m$ is found such that $f(m) \equiv 0 \ (\mod p)$. Then there is a ring homomorphism $\varphi: \mathbb{Z}[\alpha] \rightarrow \mathbb{F}_p$, where $\alpha$ is a root of $f$. Coprime integers $(a,b)$ map to the same element in $\mathbb{F}_p$ under $a - mb \mod p$ and $a-b\alpha$. Each of these pairs of coprimes are turned into a linear equation. These are solved to give the logs of the primes which are then used to solve for $x$ as in the index calculus algorithm.

As it is a derivative of index calculus, it also has subexponential time complexity.

\subsubsection{Function Field Sieve}

The Function Field Sieve \cite{Adleman1994The} is another extension of the index calculus algorithm. Discovered in 1994, it is analogous to the GNFS over function fields; $\mathbb{Z}$ is replaced with the polynomial ring $\mathbb{F}_{q[x]}$. Being a form of index calculus, it has subexponential time complexity.

\subsection{Applications of DLP}

\subsubsection{Diffie-Hellman Key Exchange}

Diffie-Hellman Key Exchange \cite{Diffie1976New} is a form of public-key encryption. It works by the following steps:

\begin{itemize}
    \item Alice and Bob both publicly agree on a prime $p$ and generator $g$.
    \item Alice picks a secret $a$ and sends Bob $A = g^a \ (mod \ p)$.
    \item Bob picks a secret $b$ and sends Alice $B = g^b \ (mod \ p)$.
    \item Both compute a shared secret $g^{(ab)} = A^b = B^a$.
\end{itemize}
An adversary knows $(p, g, A, B)$, and needs to calculate $a$ or $b$ to reveal the secret. As shown in \ref{sec:discrete-log-algorithms}, this process is intractable so it is widely used as a foundation of asymmetric encryption.

For example, let $p = 11$ and $g = 2$. Alice and Bob choose $a = 3$ and $b = 4$ respectively. As seen in Table \ref{tab:group-example}, this gives $A = 8, B = 5$. They send each other their results and calculate the shared key: 
\begin{itemize}
    \item $B^a = 5^3 \ (\mod 11) = 125 \ (\mod 11) = 4$
    \item $A^b = 8^4 \ (\mod 11) = 4096 \ (\mod 11) = 4$
\end{itemize}
Alice and Bob have successfully shared their secret without having to share the keys.

\subsubsection{ElGamal Encryption}

The ElGamal encryption system \cite{Elgamal1985A} is based on the Diffie-Hellman Key Exchange and works by the following:
\begin{itemize}
    \item Alice selects a prime $p$ and generator $g$ for the group $G = (\mathbb{Z}/p\mathbb{Z})^\times$
    \item Alice selects a random integer $a \in G$ and computes $A = g^a \ (\mod p)$; Alice's public key is $(p, g, A)$ and private key is $a$
    \item Bob selects a random integer $b$ and computes $B = g^b \ (\mod p)$
    \item Bob represents the message as an integer $m \in G$ and computes $C=m \times A^b \ (\mod p)$
    \item Bob sends the ciphertext $c = (B, C)$ to Alice
    \item Alice computes $B^{p-1-a} = B^{-a} = g^{-ab}$
    \item Alice recovers $m$ by computing $B^{-a} \times C  \mod  p = g^{-ab}mg^{ab} = m \ (\mod p)$
\end{itemize}

In this system, an eavesdropper would have access to $(g, p, A, B, C)$ and would need to calculate either $a$ or $b$ to decipher $m$.

\newpage
\section{Generative AI Critique}

\newpage
\bibliographystyle{unsrtnat}
\bibliography{report/references}

\end{document}
